\begingroup
\setlength{\parindent}{0pt}
\setlength{\parskip}{0.5\baselineskip}
\setlength{\abovedisplayskip}{8pt}
\setlength{\belowdisplayskip}{8pt}
\setlength{\abovedisplayshortskip}{8pt}
\setlength{\belowdisplayshortskip}{8pt}
\setlength{\jot}{5pt}
\clubpenalty=10000
\widowpenalty=10000
\displaywidowpenalty=10000

\part{How the Proof Works}
\label{u:guide:how-proof-works}

The three-dimensional Navier--Stokes regularity problem asks whether a smooth fluid flow can lose its smoothness in finite time. We start with smooth, divergence-free initial velocity \(u_0\), finite energy, and positive viscosity. The fluid evolves according to
\[
\partial_tu+(u\cdot\nabla)u=-\nabla p+\nu\Delta u,
\qquad \nabla\cdot u=0.
\]

\subsection*{Critical Height and Rescaling}

The problem starts with a vortex.

As it lengthens, it thins and spins faster and faster as it gets smaller.

The concern is that it could keep doing this, with each contraction taking
less time than the one before, until the core reaches zero width and the
speed becomes unbounded. Infinitely many contractions could then fit into
a finite time.
This is a Zeno cascade (``Zeno'' for infinitely many stages in finite time;
``cascade'' for the progression to smaller scales).

The equations have no minimum vortex width. To follow this possibility
down through the scales, we need to compare a vortex with a smaller,
faster version that still satisfies the equation. Reducing lengths by
\(\lambda>1\) requires multiplying velocity by \(\lambda\) and running
the evolution \(\lambda^2\) times as fast:
\[
 u_\lambda(x,t)=\lambda u(\lambda x,\lambda^2t),
 \qquad p_\lambda(x,t)=\lambda^2p(\lambda x,\lambda^2t).
\]
The chain rule gives, with every right-hand side evaluated at
\((\lambda x,\lambda^2t)\),
\[
 \begin{aligned}
 \partial_tu_\lambda&=\lambda^3\partial_tu,
 & (u_\lambda\cdot\nabla)u_\lambda&=\lambda^3(u\cdot\nabla)u,\\
 \nabla p_\lambda&=\lambda^3\nabla p,
 & \nu\Delta u_\lambda&=\lambda^3\nu\Delta u.
 \end{aligned}
\]
Every term acquires a factor \(\lambda^3\), leaving the equation unchanged
at fixed viscosity \(\nu\).
If a flow repeated this contraction, its stage durations would satisfy
\[
 \tau_j=\lambda^{-2j}\tau_0,
 \qquad \sum_{j=0}^{\infty}\tau_j
 =\frac{\tau_0}{1-\lambda^{-2}}<\infty.
\]
For the energy measurement, squared speed grows by \(\lambda^2\), but
volume shrinks by \(\lambda^3\), giving
\[
 E[u_\lambda](t)=\lambda^{-1}E[u](\lambda^2t),
 \qquad E[u]=\frac12\int_{\mathbb R^3}|u|^2\,dx.
\]
\Needspace{10\baselineskip}
The smaller, faster vortex has less energy. A bound on total energy can
miss precisely the concentration we are concerned about. The measurement
has to include how sharply the velocity changes across the shrinking core.

A Sobolev norm gives these finer changes more weight. In the Fourier
description of the velocity, a large spatial frequency \(|\xi|\) means
variation across a short distance. Weighting each frequency by
\(|\xi|^s\) gives the homogeneous Sobolev norm of order \(s\):
\[
 \|u(t)\|_{\dot H^s}^2
 =\int_{\mathbb R^3}|\xi|^{2s}|\widehat u(\xi,t)|^2\,d\xi,
 \qquad \Lambda=(-\Delta)^{1/2}.
\]
Here \(\widehat u\) is the Fourier transform of the velocity, and
\(\Lambda^s\) multiplies it by \(|\xi|^s\). Integer orders measure
ordinary spatial derivatives; fractional orders give intermediate weights
to fine structure.

The order \(s\) determines how much weight the fine scales receive.
We can find the point where that weight exactly offsets the loss of
volume. In the Fourier transform, the velocity factor
contributes \(\lambda\) and the change of spatial variables contributes
\(\lambda^{-3}\). Substituting the resulting transform into the norm
and then setting \(\xi=\lambda\eta\) gives
\[
 \begin{aligned}
 \widehat{u_\lambda}(\xi,t)
 &=\lambda^{-2}\widehat u(\xi/\lambda,\lambda^2t),\\
 \|u_\lambda(t)\|_{\dot H^s}^2
 &=\lambda^{-4}\int_{\mathbb R^3}|\xi|^{2s}
      |\widehat u(\xi/\lambda,\lambda^2t)|^2\,d\xi\\
 &=\lambda^{-4}\lambda^{2s}\lambda^3
      \|u(\lambda^2t)\|_{\dot H^s}^2
 =\lambda^{2s-1}\|u(\lambda^2t)\|_{\dot H^s}^2.
 \end{aligned}
\]
The frequency weight contributes \(\lambda^{2s}\), balancing the
remaining factor \(\lambda^{-1}\) precisely when \(2s-1=0\).
The dividing height is \(s=\tfrac12\).
For \(s<\tfrac12\), contraction reduces the norm; for \(s>\tfrac12\),
it increases the norm. At the dividing height \(s=\tfrac12\), the norm is
unchanged by the equation's rescaling. This lets us use one bound at every
vortex width, without shrinking the region of fast-moving fluid making the
bound easier to satisfy.

This is why the critical height matters: concentration cannot make this
measurement smaller simply by occupying less space. But its value can
still change as the fluid evolves. Transport and viscosity both acquired
a factor \(\lambda^3\), so the rescaling has not decided which controls
that change.

The equation gives their contributions to the critical energy directly.
Pair it with \(\Lambda u\)
and integrate the viscous term by parts. The pressure pairing vanishes
because \(\nabla\cdot\Lambda u=0\), giving
\[
 \frac12\frac{d}{dt}\|u\|_{\dot H^{1/2}}^2
 +\nu\|u\|_{\dot H^{3/2}}^2
 =-\operatorname{Re}\langle\Lambda u,(u\cdot\nabla)u\rangle.
\]
The positive term on the left measures viscous dissipation. The transport
term on the right can increase the critical energy. To compare their sizes,
use H\"older's inequality and the classical Sobolev inequality
\(\|w\|_{L^3}\le C_3\|w\|_{\dot H^{1/2}}\) on each of the three factors:
\[
 \begin{aligned}
 |\langle\Lambda u,(u\cdot\nabla)u\rangle|
 &\le\|\Lambda u\|_{L^3}\|u\|_{L^3}\|\nabla u\|_{L^3}\\
 &\le C\|u\|_{\dot H^{1/2}}\|u\|_{\dot H^{3/2}}^2,
 \qquad C=C_3^3.
 \end{aligned}
\]
Both contributions contain the \(\dot H^{3/2}\) size. Their relative
strength is governed by the critical \(\dot H^{1/2}\) norm itself.
Substitution into the balance makes this explicit:
\[
 \frac12\frac{d}{dt}\|u\|_{\dot H^{1/2}}^2
 +\bigl(\nu-C\|u\|_{\dot H^{1/2}}\bigr)
       \|u\|_{\dot H^{3/2}}^2\le0.
\]
\Needspace{7\baselineskip}
When \(C\|u\|_{\dot H^{1/2}}<\nu\), the coefficient is positive and
the critical energy decreases. If it starts below this threshold, the
decrease keeps it there. For arbitrary initial size, the coefficient need
not be positive, so this estimate cannot establish viscous control of the
nonlinear contribution. The classical difficulty is to bound the critical
norm throughout the evolution without assuming it is already small.

\subsection*{Following One Fluid Motion}

If a singularity were to actually form, the velocity would have to blow up,
which means the acceleration would have to blow up, followed by its jerk,
and so on, and so on.

\interlinepenalty=10000

The equation already contains the velocity and its time derivative.
Differentiate the whole equation in time, once and then again:
\[
 \begin{aligned}
 &\partial_tu=\nu\Delta u-(u\cdot\nabla)u-\nabla p,\\
 &\partial_t^2u=\nu\Delta\partial_tu-(\partial_tu\cdot\nabla)u
   -(u\cdot\nabla)\partial_tu-\nabla\partial_tp,\\
 &\partial_t^3u=\nu\Delta\partial_t^2u-(\partial_t^2u\cdot\nabla)u
   -2(\partial_tu\cdot\nabla)\partial_tu
   -(u\cdot\nabla)\partial_t^2u-\nabla\partial_t^2p.
 \end{aligned}
\]
The derivatives distribute across the two velocity factors, producing the
binomial coefficients. With \(V_n=\partial_t^nu\) and
\(Q_n=\partial_t^np\), the general form is
\[
 \begin{aligned}
 V_{n+1}&=\nu\Delta V_n
 -\sum_{a=0}^{n}\binom na(V_a\cdot\nabla)V_{n-a}-\nabla Q_n,\\
 -\Delta Q_n&=\partial_i\partial_j
   \sum_{a=0}^{n}\binom na V_{a,i}V_{n-a,j}.
 \end{aligned}
\]
Taking the divergence gives the second equation; decay at spatial infinity
fixes the pressure. Every \(V_n\) is a derivative of the original velocity,
and every \(Q_n\) is determined by its velocity products. All of these
derivatives belong to one fluid.

But the vortex is also narrowing across space. Each time derivative is a
spatial field, with its own spatial derivatives. Retaining those gives
\[
 J_{n,\alpha}=\partial_x^\alpha\partial_t^nu,
 \qquad \Pi_{n,\alpha}=\partial_x^\alpha\partial_t^np.
\]
Here \(\alpha\) records the spatial directions and derivative orders.
Together the fields form the mixed jet. The product rule acts in both
time and space:
\[
 J_{n+1,\alpha}=\nu\Delta J_{n,\alpha}-\nabla\Pi_{n,\alpha}
 -\sum_{a=0}^{n}\sum_{\beta\le\alpha}
   \binom na\binom\alpha\beta
   (J_{a,\beta}\cdot\nabla)J_{n-a,\alpha-\beta}.
\]
The two binomial factors count where those derivatives land. Pressure,
viscosity, and transport remain together in every equation. These time
derivatives are taken at fixed positions; along a moving fluid particle,
the derivative is \(D_t=\partial_t+u\cdot\nabla\). The transport already
present in the equations accounts for that motion.

Return to the critical energy. We left it in a balance involving the
spatial order \(3/2\). Write the energy at a general height \(s\) as
\[
 E_s=\frac12\|\Lambda^su\|_2^2.
\]
Its time derivative comes from the velocity equation. Substitution and
integration by parts give
\[
 \begin{aligned}
 E_s'
 &=\operatorname{Re}\langle\Lambda^su,\Lambda^s\partial_tu\rangle\\
 &=-\nu\|\Lambda^{s+1}u\|_2^2+\mathcal P_s\\
 &=-2\nu E_{s+1}+\mathcal P_s,
 \end{aligned}
\]
with the transport and pressure contributions collected in
\[
 \mathcal P_s=-\operatorname{Re}\langle\Lambda^su,
 \Lambda^s((u\cdot\nabla)u+\nabla p)\rangle.
\]
The viscous term contains \(E_{s+1}\): integration by parts put one
of the Laplacian's two derivatives on each velocity factor. Starting
from \(H=E_{1/2}\), differentiation of that next energy gives
\[
 \begin{aligned}
 H'&=-2\nu E_{3/2}+\mathcal P_{1/2},\\
 H''&=-2\nu E_{3/2}'+\partial_t\mathcal P_{1/2}\\
    &=4\nu^2E_{5/2}-2\nu\mathcal P_{3/2}
        +\partial_t\mathcal P_{1/2},\\
 H^{(n)}&=(-2\nu)^nE_{n+1/2}
        +\sum_{j=0}^{n-1}(-2\nu)^j
           \partial_t^{\,n-1-j}\mathcal P_{j+1/2}.
 \end{aligned}
\]
The higher we go in time derivatives, the higher we go in Sobolev space.
Every step brings another factor \(-2\nu\), with the complete transport
and pressure derivatives alongside it. All of those terms are determined
by the mixed jet of this fluid. The spatial heights run through
\[
 H^{1/2},\qquad H^{3/2},\qquad H^{5/2},\qquad H^{7/2},\qquad\ldots.
\]
We began with a vortex getting smaller and spinning faster, with its
velocity changing more sharply across the core. The higher Sobolev
orders measure precisely those sharp changes. Above height
\(k+\tfrac32\), Sobolev embedding bounds the ordinary spatial
derivatives through order \(k\):
\[
 H^s(\mathbb R^3)\hookrightarrow C_b^k(\mathbb R^3),
 \qquad s>k+\tfrac32.
\]
Any chosen derivative is covered by a sufficiently high finite level.

\phantomsection\label{u:guide:intersection}
Continuing through all finite orders reaches the full smoothness condition:
\[
 L^2\cap\bigcap_{n\ge0}\dot H^{n+1/2}(\mathbb R^3)
 =\bigcap_{m\ge0}H^m(\mathbb R^3)
 =H^\infty(\mathbb R^3)\hookrightarrow C_b^\infty(\mathbb R^3).
\]
The \(L^2\) term retains the large spatial scales; the increasing Sobolev
orders retain the small ones. This is the full spatial regularity of the
fluid.

See, the intuition here is what matters: when you follow a vortex down to
its smallest scales, highest frequencies, and fastest velocities, you are
following its successive time derivatives. For a singularity to occur,
you would expect the spatial regularity to break down as you go up in
those time derivatives. But what we see here is that viscosity relates
that faster motion to more spatial regularity, not less. This is the
opposite of what you would expect from a singularity or blowup.

\Needspace{10\baselineskip}
\vspace{2\baselineskip}
\subsection*{Proof and Application}

Establishing this regularity through a proposed singular time requires bounds
on the coupled derivatives throughout the motion.

Fix one \(u_0\in H^\infty_\sigma(\mathbb R^3)\) and one \(\nu>0\).
The subscript \(\sigma\) records incompressibility. Suppose the smooth
motion ends at a finite time \(T_*\), and write \(I=(0,T_*)\).
The initial velocity determines the pressure and, through the equation,
the subsequent changes in velocity:
\[
 \begin{aligned}
 u(t)&=u_0+\int_0^t
 \bigl[\nu\Delta u(s)-(u(s)\cdot\nabla)u(s)-\nabla p(s)\bigr]\,ds,\\
 p(t)&=R_iR_j\bigl(u_i(t)u_j(t)\bigr),
 \qquad 0\le t<T_*.
 \end{aligned}
\]
Here \(R_iR_j\) solves the pressure Poisson equation with its decaying
normalization. Substitution gives the first time derivative,
\[
 V_1(t)=\nu\Delta u(t)-(u(t)\cdot\nabla)u(t)-\nabla Q_0(t),
\]
then differentiation of the pressure product gives \(Q_1\), and the
next velocity equation gives \(V_2\). At \(t=0\), this recurrence
generates every initial derivative from \(u_0\). Throughout the motion,
it gives the mixed fields
\[
 J_{n,\alpha}(t,x)=\partial_x^\alpha V_n(t,x),
 \qquad
 \Pi_{n,\alpha}(t,x)=\partial_x^\alpha Q_n(t,x).
\]

The datum-generated pressure-complete mixed jet carries the spatial Sobolev
domains of each temporal row on \(I\). Taking every finite spatial order
gives its spatial intersection; the compatible temporal derivative identities
give the mixed intersection. The critical-height domains provide the spatial
reading during this generation. The subsequent substitution into the energy
rows displays viscosity with transport and pressure retained.

Arrange the time derivatives down the rows of a matrix and spatial orders
across its columns. A finite rectangle selects the coordinates and
relationships used in one calculation. Each entry retains all of space and
the original interval \(I\), including its adjacent temporal, viscous,
transport, and normalized-pressure coordinates.

The viscous relationship is an exact relation between entries of this
matrix. Rearranging the time equation gives
\[
 \nu\Delta V_n=V_{n+1}
 +\sum_{a=0}^{n}\binom na(V_a\cdot\nabla)V_{n-a}
 +\nabla Q_n.
\]
The next time response is tied to two more spatial derivatives, together
with the complete transport and pressure sum. The initial jet fixes where
every row begins. The construction retains these equations together on
\(I\), with pressure fixed by its velocity products, and takes finite
records of their compatible Sobolev realization. Increasing a rectangle
reads more derivatives of this motion.

A bound that deteriorates as \(t\) approaches \(T_*\) would still
leave the possible singularity unresolved. The whole-interval construction
result in Theorem~\ref{u:thm:datum-finite-blocks} establishes the paired
memberships on \(I=(0,T_*)\):
\[
 \begin{aligned}
 V_n&\in L^2(I;H^{M+1}),\\
 V_{n+1}&\in L^2(I;H^{M-1}),
 \qquad 0\le n\le N,
 \end{aligned}
\]
for every finite \(N,M\). The paired heights \(M+1\) and \(M-1\)
differ by the Laplacian's two derivatives. The rectangle
\(\cR_{N,M}[u_0;T_*]\) has finite size
\[
 \begin{aligned}
 \mathfrak R_{N,M}(u;T_*)
 &=\sum_{n=0}^{N}\Bigl(
   \|V_n\|_{L^2(I;H^{M+1})}^2\\
 &\hspace{36mm}+\|V_{n+1}\|_{L^2(I;H^{M-1})}^2\Bigr)<\infty.
 \end{aligned}
\]
Each finite depth and height has its own bound. Restricting the integral
to an earlier time can only decrease it.

Time integration alone can allow a tall, narrow spike. Here the adjacent
row also controls the rate of change. Differentiating the squared
\(H^M\) norm and splitting its spatial weights gives
\[
 \frac{d}{dt}\|V_n(t)\|_{H^M}^2
 =2\operatorname{Re}\langle
 (1-\Delta)^{(M+1)/2}V_n,
 (1-\Delta)^{(M-1)/2}V_{n+1}\rangle_{L^2}.
\]
The two factors are precisely the pair in the rectangle. Integration,
Cauchy--Schwarz, and \(2ab\le a^2+b^2\) give
\[
 \sup_{0\le t<T_*}\|V_n(t)\|_{H^M}^2
 \le\|V_n(0)\|_{H^M}^2
 +\|V_n\|_{L^2(I;H^{M+1})}^2
 +\|V_{n+1}\|_{L^2(I;H^{M-1})}^2.
\]
The initial value is fixed by \(u_0\), and the integrals remain finite
through \(T_*\). The row is bounded at every time, including the approach
to the proposed singularity. Increasing the rectangle gives this control
at each further derivative order.

Pressure follows the velocity products in its Poisson equation. A spatial
product estimate, followed by Cauchy--Schwarz in time, gives
\[
 \begin{aligned}
 \|Q_n\|_{L^1(I;H^M)}
 &\le C_M^{\rm p}\sum_{a=0}^{n}\binom na
       \|V_a\|_{L^2(I;H^{M+1})}\\
 &\hspace{35mm}\cdot\|V_{n-a}\|_{L^2(I;H^{M+1})}<\infty,
 \end{aligned}
\]
where \(C_M^{\rm p}\) is the spatial product constant. Spatial
differentiation distributes the derivatives across those factors, so a
sufficiently large rectangle covers every chosen pressure derivative too.

The rectangles agree where they overlap: an entry \((n,\alpha)\)
always means \(\partial_x^\alpha\partial_t^nu\), beginning with the
derivatives generated by \(u_0\). Enlarging the matrix preserves its
existing entries. This compatibility lets us collect their bounds without
changing the fluid being described.

Choose any time order \(r\) and spatial order \(m\). A large enough
rectangle contains \(V_0,\ldots,V_r\) in \(L^2(I;H^m)\), giving
\[
 \|u\|_{H^r(I;H^m)}^2
 =\sum_{n=0}^{r}\|V_n\|_{L^2(I;H^m)}^2<\infty.
\]
Since every finite pair \((r,m)\) is covered,
\[
 u\in\bigcap_{r,m\ge0}H^r(I;H^m(\mathbb R^3)).
\]
The full intersection means that each finite number of time and spatial
derivatives has a finite norm through \(T_*\). Every membership concerns
the fluid that began at \(u_0\).

\phantomsection\label{u:guide:mixed-jet-map}
The critical-height equations express the spatial side of this construction.
Rearranging the identity already derived isolates each higher spatial energy:
\[
 \begin{aligned}
 (-2\nu)^k E_{k+1/2}
 &=H^{(k)}
   -\sum_{j=0}^{k-1}(-2\nu)^j
       \partial_t^{\,k-1-j}\mathcal P_{j+1/2},
 \qquad k\ge1.
 \end{aligned}
\]
Positive \(\nu\) makes its coefficient nonzero. Finite selections form
the critical-height rectangles, obtained by differentiating the scalar
energy. The mixed rectangles come from differentiating the velocity and
pressure fields. Their common entries agree because both calculations
use the original equation for \(u\).

On the constructed whole-interval jet, these identities recover its
half-integer spatial energies. Kinetic energy covers low frequencies;
a weight of order \(k+\tfrac12\ge m\) covers the high frequencies
needed for any chosen integer order \(m\). Together they give every
finite spatial Sobolev order.

The equation also lets us recover the other rows from this complete
spatial information. At any chosen height \(m\), two additional spatial
orders give
\[
 \begin{aligned}
 \|(u\cdot\nabla)u\|_{L^1(I;H^m)}
 &\le C_m\|u\|_{L^2(I;H^{m+2})}^2,\\
 \|\Delta u\|_{L^1(I;H^m)}
 &\le\sqrt{T_*}\|u\|_{L^2(I;H^{m+2})},
 \end{aligned}
\]
with \(C_m\) depending on the spatial order. The pressure formula gives
an integrable pressure gradient as well. The right-hand side of the
velocity equation is now integrable through \(T_*\), so its integral
gives a continuous terminal velocity at every height. The recurrences
then give \(Q_0,V_1\), followed by \(Q_1,V_2\), and so on.
The time and spatial descriptions recover one another with viscosity,
pressure, and transport present throughout.

There is still a distinction between boundedness before \(T_*\) and
reaching a state from which the fluid can continue. The next time row
settles it. Integrating \(V_{n+1}\) between two late times gives
\[
 \begin{aligned}
 \|V_n(t)-V_n(s)\|_{H^m}
 &\le\int_s^t\|V_{n+1}(\tau)\|_{H^m}\,d\tau\\
 &\le(t-s)^{1/2}\|V_{n+1}\|_{L^2(s,t;H^m)},
 \qquad s<t<T_*.
 \end{aligned}
\]
The remaining change tends to zero. Each row reaches a limit in every
finite \(H^m\), and the limits at different heights agree. The velocity
reaches a smooth divergence-free state
\(u_*\in H^\infty_\sigma(\mathbb R^3)\). The pressure formula and
differentiated equations pass to these limits, giving its complete velocity
and pressure jet.

Local existence starts a smooth solution from \(u_*\). Its derivatives
match the incoming ones, and uniqueness joins the continuation to the
incoming motion. The fluid passes through \(T_*\). It cannot be a
finite maximal time.

\phantomsection\label{u:guide:finite-estimates}
Each finite part of the matrix also gives concrete bounds on any finite
interval \([0,T]\). For a time
derivative of order \(r\) and a spatial derivative of order \(k\), choose
\(m>k+\tfrac32\). Sobolev embedding then bounds its largest size
anywhere in the fluid:
\[
 \sup_{0\le t\le T}\|\partial_t^ru(t)\|_{W^{k,\infty}}
 \le C_{m,k}\sup_{0\le t\le T}\|\partial_t^ru(t)\|_{H^m}<\infty.
\]
Here \(C_{m,k}\) is the embedding constant. In particular, the velocity
gradient has a finite time integral. That gradient bounds how quickly
neighbouring fluid particles approach each other. For particles starting
at \(a\) and \(b\), integration along their trajectories gives
\[
 |X(a,t)-X(b,t)|\ge |a-b|
 \exp\!\left(-\int_0^t\|\nabla u(s)\|_\infty\,ds\right)>0.
\]
The vortex can stretch and narrow, but a positive material separation
retains a positive lower bound throughout each finite interval.

So, for example, choose the rectangle with \(N=0\) and \(M=1\).
It contains the velocity in \(L^2(0,T;H^2)\) and its time derivative in
\(L^2(0,T;L^2)\). Write their sizes as
\[
 \begin{aligned}
 A_T&=\left(\int_0^T\|u(t)\|_{H^2}^2\,dt\right)^{1/2},\\
 B_T&=\left(\int_0^T\|\partial_tu(t)\|_2^2\,dt\right)^{1/2}.
 \end{aligned}
\]
The adjacent-row calculation now bounds the \(H^1\) norm at every
time. Pair \(u_t\) with \((1-\Delta)u\) and integrate:
\[
 \begin{aligned}
 \|u(t)\|_{H^1}^2
 &=\|u_0\|_{H^1}^2
   +2\operatorname{Re}\int_0^t
       \langle(1-\Delta)u(s),\partial_su(s)\rangle_{L^2}\,ds\\
 &\le\|u_0\|_{H^1}^2+2A_TB_T.
 \end{aligned}
\]
These two entries give three familiar regularity estimates. First, the
\(L^3\) velocity norm is unchanged by the vortex rescaling, so it
measures concentration across scales. The spatial embedding
\[
 \|f\|_{L^3}\le\frac{\|f\|_{H^1}}{2^{5/6}\pi^{1/6}}
\]
turns the \(H^1\) bound into the Escauriaza--Seregin--\v{S}ver\'ak
endpoint estimate. The numerical constants here use the paper's unitary
Fourier normalization:
\[
 \sup_{0\le t\le T}\|u(t)\|_{L^3}
 \le \frac{\bigl(\|u_0\|_{H^1}^2+2A_TB_T\bigr)^{1/2}}
              {2^{5/6}\pi^{1/6}}.
\]

Second, the \(H^2\) entry bounds the largest speed anywhere in the fluid:
\[
 \|u(t)\|_{L^\infty}\le\frac{\|u(t)\|_{H^2}}{\sqrt{8\pi}}.
\]
Its squared time integral is already contained in \(A_T^2\), giving
the Serrin estimate,
\[
 \left(\int_0^T\|u(t)\|_{L^\infty}^2\,dt\right)^{1/2}
 \le\frac1{\sqrt{8\pi}}
       \left(\int_0^T\|u(t)\|_{H^2}^2\,dt\right)^{1/2}
 =\frac{A_T}{\sqrt{8\pi}}.
\]

\Needspace{9\baselineskip}
Third, the spatial \(H^2\hookrightarrow W^{1,3}\) estimate bounds
the velocity gradient:
\[
 \|\nabla u(t)\|_{L^3}
 \le\frac{\sqrt3\,7^{1/6}}{2^{5/3}\pi^{1/6}}\|u(t)\|_{H^2}.
\]
Its squared time integral is again bounded by \(A_T^2\), giving the
Beir\~ao da Veiga estimate~\cite{BeiraoDaVeiga1995}:
\[
 \left(\int_0^T\|\nabla u(t)\|_{L^3}^2\,dt\right)^{1/2}
 \le\frac{\sqrt3\,7^{1/6}}{2^{5/3}\pi^{1/6}}
       \left(\int_0^T\|u(t)\|_{H^2}^2\,dt\right)^{1/2}
 =\frac{\sqrt3\,7^{1/6}}{2^{5/3}\pi^{1/6}}A_T.
\]
\phantomsection\label{u:guide:how-proof-ends}
The memberships measured by \(\mathfrak R_{0,1}\) give all three estimates.
Choosing higher rectangles extends the calculation to higher derivatives.

An external force changes the acceleration of the fluid. It enters both
the velocity equation and the pressure that enforces incompressibility:
\[
 \begin{aligned}
 u_t+(u\cdot\nabla)u+\nabla p&=\nu\Delta u+f,
 &\nabla\cdot u&=0,\\
 -\Delta p&=\partial_i\partial_j(u_i u_j)-\nabla\cdot f.
 \end{aligned}
\]
Fix the initial velocity, positive viscosity, and the prescribed force
\(f\). The forced theorem takes \(f\) smooth through each finite
time, with its mixed derivatives decaying rapidly in space, uniformly on
finite intervals:
\(f\in C^\infty([0,\infty);\mathcal S(\mathbb R^3;\mathbb R^3))\).
A smooth force with compact support in space and time belongs to this
class. Its size enters the bounds. Writing \(F_n=\partial_t^nf\),
the differentiated equations become
\[
 \begin{aligned}
 V_{n+1}&=\nu\Delta V_n
 -\sum_{a=0}^{n}\binom na(V_a\cdot\nabla)V_{n-a}
 -\nabla Q_n+F_n,\\
 -\Delta Q_n&=\partial_i\partial_j
   \sum_{a=0}^{n}\binom na V_{a,i}V_{n-a,j}-\nabla\cdot F_n.
 \end{aligned}
\]
Spatial differentiation adds \(\partial_x^\alpha F_n\) to each mixed
row. The initial recurrence uses \(F_n(0)\) alongside \(u_0\), and
the pressure responds to the force through its divergence.

In the energy balance, the additional term is the work done by the force:
\[
 E_s'=-2\nu E_{s+1}+\mathcal P_s+W_s,
 \qquad
 W_s=\operatorname{Re}\langle\Lambda^su,\Lambda^sf\rangle.
\]
The viscous coefficient stays \(-2\nu\). Repeated differentiation
now includes the derivatives of the force work:
\[
 H^{(k)}=(-2\nu)^k E_{k+1/2}
 +\sum_{j=0}^{k-1}(-2\nu)^j
  \partial_t^{\,k-1-j}(\mathcal P_{j+1/2}+W_{j+1/2}).
\]
Every derivative of \(W_s\) pairs velocity derivatives with known
derivatives of \(f\). Their contribution to a selected row can be
bounded directly. Split the spatial weights between \(V_n\) and
\(F_n\), then apply Cauchy--Schwarz and Young's inequality:
\[
 2\bigl|\langle V_n,F_n\rangle_{H^M}\bigr|
 \le\nu\|V_n\|_{H^{M+1}}^2
      +\nu^{-1}\|F_n\|_{H^{M-1}}^2.
\]
The first term is absorbed at the spatial order of the viscous estimate.
The second has a finite time integral because the force is prescribed
and smooth through \(T_*\). Its pressure contribution is controlled by
the Poisson formula. Retaining these source terms, the forced construction
supplies the paired whole-interval bounds
\[
 \mathfrak R^f_{N,M}(u;T_*)
 =\sum_{n=0}^{N}\left(
 \|V_n\|_{L^2(I;H^{M+1})}^2
 +\|V_{n+1}\|_{L^2(I;H^{M-1})}^2\right)<\infty
\]
for every finite \(N,M\). The bounds now depend on the force as well
as \(u_0\) and \(\nu\). The adjacent-row calculation again gives
bounds at every time, and compatibility collects the finite orders into
\[
 u\in\bigcap_{r,m\ge0}H^r(I;H^m(\mathbb R^3)).
\]

Integrating the next time derivative makes the remaining change tend to
zero at \(T_*\). The fluid reaches a smooth state \(u_*\), with
pressure and higher derivatives fixed by the forced recurrence and the
values \(F_n(T_*)\). Local existence continues from that state with
the prescribed \(f(t)\) at its original time. The derivatives agree
across \(T_*\), and uniqueness joins the incoming and outgoing motion.
The proposed singular time is again an ordinary time of the flow.

The bounds reach each term of the momentum equation individually:
\[
 \sup_{0\le t\le T}\left(
 \|u_t\|_{H^m}+\|(u\cdot\nabla)u\|_{H^m}
 +\|\nabla p\|_{H^m}+\nu\|\Delta u\|_{H^m}
 \right)<\infty,
 \qquad m<\infty.
\]
The time derivative is a row of the matrix; the Laplacian uses two more
spatial orders; transport is a product of controlled derivatives; pressure
is determined by the velocity and force. Sobolev embedding bounds their
pointwise sizes too. Thus each term stays finite, even when the terms
partly cancel in the equation.

OpenAI's proposed blowup starts from rest, uses a smooth force with compact support,
and concentrates increasing velocity into a shrinking region while keeping
kinetic energy finite.\footnote{OpenAI,
\href{https://cdn.openai.com/pdf/32d9f210-8b73-45e0-91bc-82a30aef8a9a/navier-stokes.pdf}{\emph{Finite Time Blowup for Navier--Stokes}},
Theorem~1.1 and Section~10, especially (10.20)--(10.21).}
This returns to the opening problem: total energy can remain finite while
velocity concentrates into a smaller region. Write \(v\) for the
proposed velocity and put its proposed singular time at \(t=1\).
With \(\tau=1-t\), its claimed growth in (10.21) implies
\[
 \|v(1-\tau)\|_{L^\infty}\ge c\,\tau^{-A}\longrightarrow\infty,
 \qquad c>0,\quad A>0.
\]
The relevant bound is on the largest speed. If the proposed force has
its stated smooth extension through \(t=1\), denote it by
\(f_{\rm OA}\). The forced theorem gives a smooth solution \((U,P)\)
from rest for that force and viscosity, through \(t=1\).

Before \(t=1\), both \(v\) and \(U\) solve the prescribed initial
value problem. Their difference \(z=v-U\) starts at zero. Subtracting
the equations cancels the force; pairing with \(z\) and using
incompressibility gives
\[
 \frac12\frac{d}{dt}\|z\|_2^2+\nu\|\nabla z\|_2^2
 \le\|\nabla U\|_\infty\|z\|_2^2,
 \qquad z(0)=0.
\]
The coefficient has a finite time integral, so Gronwall's inequality
keeps the difference zero. On every closed interval ending before \(1\),
\(v=U\).

The rectangle \(N=0,M=2\) contains \(U\) at spatial height \(3\)
and \(U_t\) at height \(1\). Their adjacent-row trace bounds
\(H^2\) at every time, which bounds the speed. With zero initial velocity,
\[
 \begin{aligned}
 \mathfrak R^f_{0,2}[U]
 &=\|U\|_{L^2(0,1;H^3)}^2
   +\|U_t\|_{L^2(0,1;H^1)}^2<\infty,\\
 \sup_{0\le t<1}\|v(t)\|_{L^\infty}^2
 &=\sup_{0\le t<1}\|U(t)\|_{L^\infty}^2
 \le\frac{\mathfrak R^f_{0,2}[U]}{8\pi}<\infty.
 \end{aligned}
\]
The proposed divergence contradicts this finite bound. Under the stated
smooth-force hypothesis, the flow continues smoothly through \(t=1\);
the divergent field cannot solve that initial value problem.

The vortex can stretch, narrow, and gain energy from the force. Its velocity,
pressure, and derivatives remain coupled throughout that motion. The
finite bounds obtained from their viscous relationship survive down through
the scales, and the fluid reaches a smooth state from which it continues.

\endgroup
